\section{Top K Frequent Elements}
\label{sec:heap-top-k-frequent}

\problemstatement{Given an array $\textit{nums}$ and an integer $k$,
return the $k$ elements that appear \textbf{most frequently} in the
array.}

\begin{patternbox}
\emph{``Top $k$ frequent''} is the size-$k$ heap technique from
Section~\ref{sec:heap-kth-largest}, applied not to the raw values
themselves but to a \textbf{derived key} -- here, frequency, computed once
via a hash map before the heap ever comes into play. Whenever a problem
asks for the top/bottom $k$ items ranked by something other than their
own raw value (frequency, distance, a custom score), first compute that
key for every item, then apply the familiar size-$k$ heap exactly as if
the key were the value.
\end{patternbox}

\subsection*{Understanding the problem carefully}

First, count the frequency of every distinct element using a hash map --
this converts the problem from ``top $k$ frequent elements'' into ``top
$k$ elements by a number we now have for each of them,'' which is exactly
the shape of Section~\ref{sec:heap-kth-largest}, just with $(\textit{key} =
\textit{frequency}, \textit{value} = \textit{element})$ pairs instead of
raw numbers.

\begin{ideabox}[Reuse, with a Derived Key]
Build a frequency map. Maintain a min-heap of size $k$, ordered by
frequency, holding $(\textit{frequency}, \textit{element})$ pairs.
Process each distinct element: push it; if the heap exceeds size $k$, pop
the pair with the smallest frequency (the current weakest member of the
top-$k$-frequent set). At the end, the heap holds exactly the $k$ most
frequent elements.
\end{ideabox}

\subsection*{Why the size-$k$ heap argument transfers unchanged}

The correctness argument from Section~\ref{sec:heap-kth-largest} never
actually depended on the heap being keyed by the element's own value --
it only depended on the key being a single comparable number that we want
to rank by. Replacing ``value'' with ``frequency of that value'' changes
nothing about why the size-$k$ min-heap correctly converges to the top-$k$
set: the same induction (the heap always holds the best $k$ keys seen so
far, by the same push-then-evict-the-worst logic) applies verbatim.

\subsection*{Algorithm}

\begin{enumerate}
  \item Build a frequency map of $\textit{nums}$.
  \item For each distinct $(\textit{element}, \textit{freq})$ pair:
  \begin{enumerate}
    \item Push $(\textit{freq}, \textit{element})$ onto a min-heap
          (ordered by \textit{freq}).
    \item If the heap's size exceeds $k$: pop the minimum.
  \end{enumerate}
  \item Return the elements remaining in the heap.
\end{enumerate}

\begin{algorithm}[H]
\caption{Top K Frequent Elements}
\begin{algorithmic}[1]
\Procedure{TopKFrequent}{\textit{nums}, $k$}
  \State build frequency map \textit{freq} of \textit{nums}
  \State $\textit{heap} \gets$ empty min-heap (by frequency)
  \For{each $(\textit{elem}, f)$ in \textit{freq}}
    \State push $(f, \textit{elem})$ onto \textit{heap}
    \If{$|\textit{heap}| > k$} pop the minimum \EndIf
  \EndFor
  \State \Return all elements remaining in \textit{heap}
\EndProcedure
\end{algorithmic}
\end{algorithm}

\subsection*{Dry run}

\dryrun{Take $\textit{nums} = [1,1,1,2,2,3]$, $k = 2$.}

Frequency map: $1{:}3, 2{:}2, 3{:}1$.

\begin{itemize}
  \item Push $(3,1)$: heap $=\{(3,1)\}$.
  \item Push $(2,2)$: heap $=\{(2,2),(3,1)\}$. Size $2 \le 2$.
  \item Push $(1,3)$: heap $=\{(1,3),(2,2),(3,1)\}$. Size $3>2$: pop min
        $(1,3)$. Heap $=\{(2,2),(3,1)\}$.
\end{itemize}

The final heap holds elements $\{2, 1\}$, the two most frequent values
(frequencies $2$ and $3$ respectively).

\subsection*{C++ implementation}

\begin{lstlisting}
#include <bits/stdc++.h>
using namespace std;

vector<int> topKFrequent(vector<int>& nums, int k) {
    unordered_map<int,int> freq;
    for (int x : nums) freq[x]++;

    priority_queue<pair<int,int>, vector<pair<int,int>>,
                   greater<pair<int,int>>> minHeap; // (freq, element)

    for (auto& [elem, f] : freq) {
        minHeap.push({f, elem});
        if ((int)minHeap.size() > k) minHeap.pop();
    }

    vector<int> result;
    while (!minHeap.empty()) {
        result.push_back(minHeap.top().second);
        minHeap.pop();
    }
    return result;
}

int main() {
    vector<int> nums = {1, 1, 1, 2, 2, 3};
    for (int x : topKFrequent(nums, 2)) cout << x << " ";
    cout << endl;
    return 0;
}
\end{lstlisting}

\begin{complexitybox}
\textbf{Time Complexity.} Building the frequency map costs $O(n)$.
Processing each of the $d \le n$ distinct elements costs $O(\log k)$,
giving an overall time complexity of
\[
  O(n + d \log k).
\]
\textbf{Space Complexity.} $O(d)$ for the frequency map plus $O(k)$ for
the heap.
\end{complexitybox}

\begin{takeawaybox}
This closing problem of the chapter is a deliberate reminder: nearly
every heap problem you will face is one of a small number of templates --
size-$k$ heap, merge-$k$-sources, two-balanced-heaps, or
always-take-the-best-remaining-greedy -- applied to whatever key the
problem actually cares about ranking by. Compute that key first (here,
frequency, via a hash map), and the heap logic that follows is almost
always one you have already written in this chapter.
\end{takeawaybox}
